\section{Spectral Transfer and Dispersive Estimates}
\label{sec:dispersive_estimates}

This appendix does not prove a spectral transfer theorem. It formulates a mathematically explicit transfer problem and isolates sufficient intermediate estimates that would make such a theorem plausible.

\subsection{The transfer problem}

Let $P$ denote the linearized hyperbolic operator governing a chosen field variable on a black-hole background, schematically
\[
    P = -\partial_t^2 + \mathcal A + \mathcal E,
\]
where $\mathcal A$ is an elliptic operator on spatial slices and $\mathcal E$ contains lower-order geometric error terms. The guiding question is whether a horizon-local elliptic estimate for an operator related to MOTS/Jang geometry can be converted into spacetime decay for $P$.

\begin{proposition}[Conditional spectral-transfer principle]\label{prop:SpectralTransfer}
Suppose that the MOTS stability operator $L_\Sigma = -\Delta + 2X \cdot \nabla + Q$ exhibits a uniform spectral gap, meaning the principal eigenvalue $\lambda_1 > 0$ and the non-real spectral values relevant to the model are separated from the low-frequency axis. Assume in addition that this horizon-local coercivity transfers to the spatial part $\mathcal A$ of the spacetime operator through a valid mode decomposition and resolvent gluing argument. Then the following low-frequency resolvent bound would hold for the linearized spacetime operator $P$:
\[
    \|(P - \omega^2)^{-1}\|_{L^2_{comp} \to L^2_{loc}} \le C
\]
for all $\omega \in \mathbb{R}$ near zero.
\end{proposition}

\begin{proof}
Let $u$ be a solution to the localized elliptic equation $(\mathcal A - \omega^2)u = f$. By separating variables in a collar neighborhood of the horizon, we expand $u = \sum u_n(r)Y_n$, where $Y_n$ are the eigenmodes of the operator $\mathcal A_\Sigma$ induced by the Jang metric.
Under the assumed transfer from $L_\Sigma$ to $\mathcal A_\Sigma$, the angular operator is coercive, yielding
\[
    \langle \mathcal A_\Sigma u, u\rangle \ge \eta^2 \|u\|^2_{L^2(\Sigma)} .
\]
Consequently, for small $\omega$, the radial equations have no turning points and uniformly exhibit exponential decay into the trapped region. 
This provides the essential bound (S1), which, combined with standard high-frequency Morawetz multipliers, yields integrated local energy decay.
\end{proof}

\subsection{A sufficient package of estimates}

\begin{proposition}[Conditional spectral-transfer package]\label{prop:ConditionalSpectralTransfer}
Suppose there exist constants $c_0,C_0>0$ and a decomposition of spatial frequencies into low and high regimes such that:
\begin{enumerate}
    \item[(S1)] \textbf{Horizon coercivity.} The low-frequency component $u_{\mathrm{low}}$ satisfies
    \[
        c_0\|u_{\mathrm{low}}\|_{H^1}^2
        \le
        \langle \mathcal A u_{\mathrm{low}},u_{\mathrm{low}}\rangle
        + C_0\|u_{\mathrm{low}}\|_{L^2(K)}^2
    \]
    for some fixed compact region $K$ containing the trapped set.
    \item[(S2)] \textbf{Resolvent gluing.} For every real $\omega$ with $|\omega|\le \omega_0$, solutions of
    \[
        (\mathcal A - \omega^2 + \mathcal E_\omega)u=f
    \]
    satisfy a resolvent estimate controlled by the coercive norm from (S1).
    \item[(S3)] \textbf{High-frequency Morawetz estimate.} Frequencies $|\omega|\ge \omega_0$ obey the usual microlocal Morawetz/positive-commutator estimate.
\end{enumerate}
Then solutions of $P\psi=0$ satisfy an integrated local energy decay estimate on every finite-energy time slab.
\end{proposition}

\begin{proof}
Estimate (S2) controls the low-frequency resolvent, which is the only regime not already covered by high-frequency propagation methods. Estimate (S3) handles the high-frequency part. A standard contour argument or Fourier synthesis then combines the two bounds into an integrated local energy decay estimate for $P$. The substantive new content would therefore lie in verifying (S1)--(S2) from geometric input, not in the final functional-analytic assembly.
\end{proof}

\subsection{Where Theorem~\ref{thm:AdjointConeEuler} may matter}

The Penrose problem suggests one possible mechanism for producing (S1): if localized sign information on horizon sections can be extracted from constrained variational data, then the resulting coercive form may help rule out zero-frequency obstructions concentrated near the horizon. In that sense, Theorem~\ref{thm:AdjointConeEuler} is not itself a dispersive statement, but it points toward the sort of adjoint coercivity one would want in a future spectral-transfer theory.

\begin{remark}[Roadmap rather than theorem]
This section makes no claim that the linearized Einstein or Teukolsky operators have been reduced to the above model, nor that the needed coercive forms have been established. The purpose is only to record a clean reduction: the missing bridge is a low-frequency resolvent theorem driven by horizon-local geometry.
\end{remark}
